% ch-background.tex — Background and prerequisites

\chapter{Background}\label{cha:background}

This chapter collects the background material needed in subsequent chapters.

\section{Notation and Conventions}\label{sec:notation}

Throughout this monograph, we write $\NN$, $\ZZ$, $\QQ$, $\RR$, and $\CC$
for the natural numbers, integers, rationals, reals, and complex numbers,
respectively.

\begin{notation}\label{notn:sequences}
  A sequence in a set $X$ is a function $x \colon \NN \to X$.
  We write $(x_n)_{n \in \NN}$ or simply $(x_n)$ when the index set is clear.
\end{notation}

\section{Prerequisites from Analysis}\label{sec:prerequisites-analysis}

We recall the basic definitions and results from analysis that will be used later.

\begin{definition}\label{def:metric-space}
  A \emph{metric space} is a pair $(X, d)$ where $X$ is a set and
  $d \colon X \times X \to [0, \infty)$ is a function satisfying:
  \begin{enumerate}
    \item $d(x, y) = 0$ if and only if $x = y$;
    \item $d(x, y) = d(y, x)$ for all $x, y \in X$;
    \item $d(x, z) \le d(x, y) + d(y, z)$ for all $x, y, z \in X$.
  \end{enumerate}
\end{definition}

\begin{example}\label{exa:euclidean-metric}
  The real line $\RR$ with the absolute value metric $d(x, y) = \abs{x - y}$
  is a metric space.
\end{example}

\begin{lemma}\label{lem:triangle-inequality}
  In any metric space $(X, d)$, for all $x_1, \ldots, x_n \in X$,
  \begin{equation}\label{eq:generalized-triangle}
    d(x_1, x_n) \le \sum_{i=1}^{n-1} d(x_i, x_{i+1}).
  \end{equation}
\end{lemma}

\begin{proof}
  By induction on $n$, using the triangle inequality.
\end{proof}

%%% Local Variables:
%%% mode: LaTeX
%%% TeX-master: "monograph-template"
%%% End:
